\section{Introduction}

Many production uses of language models are classification problems in disguise: routing a support ticket, choosing an
intent, deciding whether a message needs a human. In 2026 several providers started offering \emph{decision models}
for exactly this job: endpoints that receive a text and a list of options and return a probability for each option
instead of generated text. TypeSafe AI's Jev popularised the category \citep{venturebeat2026jev}; OpenAI Decisions
\citep{openai2026decisions} and Base Labs' Elephant-1 and Mammoth-1 \citep{baselabs2026decision} followed. They are
marketed as cheap and accurate without any task-specific training. Jev itself is not part of this study: at the time of
the experiments access was limited to a waiting list.

A practitioner's natural objection is that a small open model, fine-tuned on the task and run on hardware one already
owns, should be cheaper and at least as accurate. Are decision models a meaningful new tool, or can small task-specific
classifiers, possibly run locally, achieve the same result? This paper tests that objection empirically, under explicit and
modest conditions: two public intent benchmarks (Banking77 in English, MASSIVE in Italian), a synthetic Italian
customer-service task with a lexically adversarial test condition and a phenomenon-level robustness set, consumer
hardware (a 6\,GB laptop GPU and a free 16\,GB cloud GPU), and a capped API budget.

We ask: \emph{under which conditions do zero-shot decision models outperform compact classifiers adapted to the task,
and how does the comparison change with the label budget, the category descriptions and a shift of the test
distribution?} Concretely:
\begin{itemize}
  \item[RQ1] How many labelled examples per class does a compact supervised classifier need to match zero-shot decision models,
  and does the answer depend on the training recipe?
  \item[RQ2] How much do category descriptions change zero-shot performance when everything else is held fixed?
  \item[RQ3] How much of a decision model's advantage remains when it is compared with its own base checkpoint used zero-shot?
  \item[RQ4] How do the systems behave on a lexical stress test and on targeted phenomena, including inputs that should be
  routed to a human because they are ambiguous or outside the taxonomy?
  \item[RQ5] What do calibration, latency and cost imply for deployment?
\end{itemize}

Figure~\ref{fig:design} summarises the design. Our contributions are:
\begin{itemize}
  \item an independent, paired comparison of three commercial decision models with six families of compact classifiers on two
  public benchmarks and a new synthetic Italian task, across three label budgets and three seeds;
  \item a controlled measurement of the effect of category descriptions, option order and in-context examples on the same messages,
  with a label-free comparison (label names only) separated from descriptions written from training data;
  \item a comparison of Elephant-1 with its own base checkpoint used zero-shot, which shows how much smaller the observed gap is
  against the right progenitor (not a causal ablation);
  \item a diagnosis of the few-shot failure of decoder classifiers, which persists with both recipes tested, and a sensitivity
  analysis showing that with the full training set the outcome on Banking77 depends on the recipe;
  \item a lexical stress set and a phenomenon-level robustness set with out-of-taxonomy and ambiguous messages, and an analysis of
  calibration and routing to human review;
  \item code, frozen protocol, per-example predictions and synthetic data, sufficient to recompute every reported number.
\end{itemize}

\begin{figure}[t]
\centering
\begin{tikzpicture}[
  font=\footnotesize, execute at begin node={\hyphenpenalty=10000\exhyphenpenalty=10000},
  box/.style={draw, rounded corners=2pt, align=center, inner sep=4pt, text width=4.4cm, minimum height=1.25cm},
  cond/.style={align=center, text width=4.4cm, font=\scriptsize, text=black!70, minimum height=1.05cm},
  data/.style={draw, fill=black!4, rounded corners=2pt, align=center, inner sep=4pt, text width=14.6cm},
  arr/.style={-{Stealth[length=2mm]}, thick}
]
\node[data] (tasks) at (0,0) {\textbf{Same test messages:} Banking77 (EN) $\cdot$ MASSIVE-it (IT) $\cdot$ Ottavo (IT, synthetic)
  $\cdot$ Ottavo lexical stress $\cdot$ Ottavo-R phenomena};

\node[box, fill=blue!6, anchor=north] (gen) at (0,-0.9) {\textbf{Generalist LLMs}\\ Qwen3-8B, Qwen3.5-9B\\ (base of Elephant-1)};
\node[box, fill=blue!12, left=0.3cm of gen] (dm) {\textbf{Decision models}\\ OpenAI Decisions, Elephant-1,\\ Mammoth-1 (API)};
\node[box, fill=orange!10, right=0.3cm of gen] (sup) {\textbf{Compact classifiers}\\ TF-IDF, SetFit, mmBERT,\\ Qwen3 0.6B / 1.7B / 8B};

\node[cond, below=0.1cm of dm] (cdm) {names only (N): no task labels;\\ names + descriptions (D)$^*$;\\ option order; examples in the prompt};
\node[cond, below=0.1cm of gen] (cgen) {Ottavo only, no task labels;\\ descriptions (D) = specification;\\ scoring of option letters};
\node[cond, below=0.1cm of sup] (csup) {trained on 5 or 20 examples per class\\ or on all training data; 3 seeds;\\ recipe chosen by CV inside the budget};

\node[data, anchor=north] (out) at ([yshift=-0.45cm] cgen.south |- csup.south) {\textbf{Paired comparison on identical messages:} macro-F1 with
  bootstrap CIs $\cdot$ robustness by phenomenon $\cdot$ calibration and routing to human review $\cdot$ latency and cost};

\foreach \n in {dm, gen, sup} { \draw[arr] (tasks.south -| \n.north) -- (\n.north); }
\foreach \c in {cdm, cgen, csup} { \draw[arr] (\c.south) -- (\c.south |- out.north); }
\end{tikzpicture}
\caption{Study design. The same test messages are classified by decision models used zero-shot, by generalist LLMs used the same
way, and by compact classifiers trained on increasing numbers of labelled examples. Zero-shot systems see no labelled example in the
request; with label names (N) they use no task label at all. $^*$On Banking77 and MASSIVE-it the descriptions were written from
labelled training examples.}
\label{fig:design}
\end{figure}


This is an independent technical report, not a claim of novelty: we did not find a published independent comparison of these
specific systems, but we do not claim to be the first study of zero-shot versus fine-tuned classification. Our contribution is a
pre-registered-in-repository protocol, paired comparisons on identical examples, per-example predictions released for
re-analysis, and an explicit account of what could not be tested.
